\input{algebra1.tex}
\title[Algebra 1.]{Polinomok Euklideszi-algoritmusa}
\subtitle{}
\date[3. hét]{2026.~szeptember 28-án}
\begin{document}
\openingframes{Tartalom}
\section{Polinomok gyökeiről}
\frame{
A  maradékos osztás módszerének másik fontos következménye a polinomok gyökszámáról szóló állítás.
Azt gondoljuk meg,
hogy a gyöktényező a polinomból mindig kiemelhető,
emiatt egy akármilyen test feletti $n$-ed fokú polinom gyökeinek száma $n$-nél nagyobb nem lehet.
\begin{proposition}
	Legyen $p\in\mathbb{F}\left[ t \right]$ egy legalább első fokú polinom,
	és $t_0$ annak egy gyöke, azaz $p\left( t_0 \right)=0$.
	Ekkor létezik $h\in\mathbb{F}\left[ t \right]$ polinom,
	amelyre
	\[
		\deg h= \deg p -1
        \quad\text{ és }\quad
        p\left( t \right)=\left( t-t_0 \right)h\left( t \right).\qedhere
	\]
\end{proposition}
}
%\begin{proof}
%	Maradékos osztással $p$-re és az elsőfokú $t-t_0$ polinomra
%	\[
%		p\left( t \right)=h\left( t \right)\left( t-t_0 \right)+r\left( t \right),
%		\text{ ahol }
%		\deg r<1.
%	\]
%	No de, $t_0$ egy gyök, tehát $0=p\left( t_0 \right)=r\left( t_0 \right)$.
%	Ez azt jelenti, hogy $\deg r=-\infty$, ami éppen az állítás.
%\end{proof}
\frame{
\begin{definition}[gyök multiplicitása]\index{gyökök multiplicitása}
	Legyen $t_0$ gyöke a $p\left( t \right)$ polinomnak.
	Azt mondjuk, hogy a $k$ pozitív egész e $t_0$ gyök \emph{multiplicitása},
	ha van olyan $h\left( t \right)$ polinom, hogy
	\[
		p\left( t \right)=\left( t-t_0 \right)^kh\left( t \right),
        \text{ de }
        h\left( t_0 \right)\neq 0.
    \]
	Néha azt is mondjuk, hogy $t_0$ egy $k$-szoros gyöke $p$-nek.
\end{definition}
Világos, hogy a gyöktényező kiemelhetősége miatt egy nem zérus polinom minden gyökének van multiplicitása, és az legalább $1$.
\begin{proposition}
	Legyen(ek) a $p\in\mathbb{F}\left[ t \right]$ nem zérus polinom különböző gyökei $t_1,\ldots,t_k$,
	és ezen gyökök multiplicitásai rendre $m_1,\ldots,m_k$.
	Ekkor 
    $$m_1+\ldots+m_k\leq\deg p.$$
\end{proposition}
%\begin{proof}
%	A test nullosztó mentessége és a gyöktényező kiemelhetősége miatt
%	\[
%		p\left( t \right)=
%		\left( t-t_1 \right)^{m_1}\cdot\left( t-t_2 \right)^{m_2}\ldots\left( t-t_k \right)^{m_k}\cdot
%		h\left( t \right),
%	\]
%	ahol $h$ olyan nem zérus polinom, amelynek már nincsen gyöke.
%	A fokszámok összehasonlításából kapjuk, hogy
%	$m_1+\ldots+m_k\leq m_1+\ldots+m_k+\deg h=\deg p$.
%\end{proof}
}
\frame{
A fenti gondolat szerint, ha egy legfeljebb $n$-ed fokú polinomnak $n+1$ különböző gyöke van,
akkor csak úgy lehetséges, ha a polinom minden együtthatója nulla.
Ezt úgy is szoktuk fogalmazni,
hogy egy legfeljebb $n$-ed fokú polinomot $n+1$ helyettesítési értéke már egyértelműen meghatározza:
\begin{proposition}
	Tegyük fel, hogy a $p,q\in\mathbb{F}\left[ t \right]$ polinomok legfeljebb $n$-ed fokúak, ahol $n$ egy nemnegatív egész,
	és tegyük fel,
	hogy létezik $n+1$ különböző $t_0,t_1,\ldots,t_n$ pont a testben,
	amelyekre $p\left( t_j \right)=q\left( t_j \right)$ minden $j=0,\ldots,n$.
	Ekkor $p\left( t \right)=q\left( t \right)$,
	azaz $p$ és $q$ együtthatói azonosak.
    \label{pr:polinomfv}
\end{proposition}
%\begin{proof}
%	Legyen $h=p-q$.
%	Világos, hogy $\deg h\leq n$ és $h$-nak van $n+1$ különböző gyöke.
%	Az előző állítás szerint ez csak a $h=0$ polinomra igaz, ami azt jelenti,
%	hogy $p$ és $q$ együtthatói azonosak.
%\end{proof}
}
\frame{
A maradékos osztás tételének szép következménye tehát,
hogy ha $\mathbb{F}$ egy nem véges test, és a $p,q\in\mathbb{F}\left[ t \right]$ polinomok,
akkor $p$-nek és $q$-nak pontosan akkor azonosak az együtthatói, ha
\(
p\left( t \right)=q\left( t \right)
\)
fennáll minden $t\in\mathbb{F}$ mellett.

Itt fontos, hogy $\mathbb{F}$ nem véges test,
hiszen például ha $\mathbb{F}=\left\{ 0,1 \right\}$ a két elemű test,
akkor a $p\left( t \right)=t^2+t$ polinomra minden $t\in\left\{ 0,1 \right\}$ mellett $p\left( t \right)=0$,
de a polinom együtthatói rendre a $\left\{ 0,1,1 \right\}$ számok a testből,
tehát ez nem a összeadásra nézve neutrális eleme az $\mathbb{F}\left[ t \right]$ polinomgyűrűnek.
}

\section{Polinomok legnagyobb közös osztója és legkisebb közös többszöröse}
\frame{
\begin{definition}
	Legyenek most $p_1,\ldots,p_k$ polinomok.
	\begin{enumerate}
		\item A $d$ polinom a \emph{legnagyobb közös osztója}\index{legnagyobb közös osztó} az adott polinomoknak,
		      ha
		      \begin{enumerate}
			      \item $d|p_j$ minden $j=1,\ldots,k$-ra,
			      \item ha $d_1|p_j$ minden $j=1,\ldots,k$ mellett akkor $d_1|d$ is fennáll,
			      \item $d$ normált.
		      \end{enumerate}
		      A $p_1,\ldots,p_k$ polinomokat \emph{relatív prímeknek}\index{relatív prím polinomok} nevezzük,
		      ha közös osztójuk csak a konstans polinomok,
		      azaz a $d\left( t \right)=1$ a legnagyobb közös osztó.
		\item A $d$ polinom a \emph{legkisebb közös többszöröse}\index{legkisebb közös többszörös} az adott polinomoknak,
		      ha
		      \begin{enumerate}
			      \item $p_j|d$ minden $j=1,\ldots,k$-ra,
			      \item ha $p_j|d_1$ minden $j=1,\ldots,k$ mellett akkor $d|d_1$ is fennáll.
			      \item $d$ normált.\qedhere
		      \end{enumerate}
	\end{enumerate}
\end{definition}
Persze az első kérdés, hogy van-e a polinomoknak legnagyobb közös osztója vagy legkisebb közös többszöröse, és hány ilyen van?
}
\frame{
\begin{proposition}
	Bármely $p_1,\ldots,p_r\in\mathbb{F}\left[ t \right]$ nem zérus polinomoknak
	létezik egyetlen legkisebb közös többszöröse.

	Jelesül, a legkisebb közös többszörös, a
	\[
	\cap_{j=1}^rJ\left( p_j \right)
    \]
    ideálnak, mint főideálnak az egyetlen normált generátora.
    \label{pr:lkkt}%
\end{proposition}
%\begin{proof}
%	Világos, hogy ideálok metszete is ideál, emiatt
%	\(
%	\cap_{j=1}^rJ\left( p_j \right)
%	\)
%	is ideál $\mathbb{F}\left[ t \right]$ gyűrűben.
%	De itt minden ideál főideál, létezik tehát $d\in\mathbb{F}\left[ t \right]$ normált polinom, amelyre
%	\[
%		J\left( d \right)
%		=
%		\cap_{j=1}^rJ\left( p_j \right)
%	\]
%	Világos, hogy $d\in J\left( p_j \right)$ minden $j$-re,
%	ergo $d$ többszöröse minden $p_j$-nek.
%	Ha $p_j|d_1$ fennáll, minden $j$-re
%	az azt jelenti, hogy $d_1\in J\left( p_j \right)$, minden $j$-re, azaz
%	\(
%	d_1
%	\in
%	\cap_{j=1}^rJ\left( p_j \right)
%	=
%	J\left( d \right),
%	\)
%	tehát $d|d_1$ valóban fennáll.
%
%	Ha $d$ mellett $g$ is legkisebb közös többszörös, akkor $d|g$ és $g|d$ szerint $g$ és $d$ foka azonos,
%	így csak egymás konstans szorosai lehetnek, de mivel mindketten normáltak, ezért e konstans csak 1 lehet.
%\end{proof}
}
\frame{
\begin{proposition}
	Bármely $p_1,\ldots,p_r\in\mathbb{F}\left[ t \right]$ nem zérus polinomoknak
	létezik egyetlen legnagyobb közös osztója.

	Jelesül,
	a legnagyobb közös osztó a $J\left( p_1,\ldots,p_r \right)$ ideálnak mint főideálnak az egyetlen generáló eleme.
	Ezért a $d$ legnagyobb közös osztó kifejezhető
	\[
		d\left( t \right)=f_1\left( t \right)p_1\left( t \right)+
		\ldots+
		f_r\left( t \right)p_r\left( t \right)
	\]
	alakban,
	valamely $f_1,\ldots,f_r\in\mathbb{F}\left[ t \right]$ polinomok segítségével.
\end{proposition}
%\begin{proof}
%	Láttuk, hogy létezik egyetlen normált $d$ polinom, amelyre $J\left( d \right)=J\left( p_1,\ldots,p_r \right)$.
%	Világos, hogy $d|p_j$ minden $j=1,\ldots,r$ és $d\in J\left( p_1,\ldots,p_r \right)$,
%	azaz
%	\[
%		d=f_1p_1+\ldots+f_rp_r
%	\]
%	valamely $f_1,\ldots,f_r$ polinomokra.
%	Ha valamely $d_1$ polinomra $d_1|p_j$ minden $j=1,\ldots,r$ mellett,
%	akkor a fenti azonosság szerint $d_1|d$ is fennáll.
%
%	Az egyértelműség mint a legkisebb közös többszörösnél.
%\end{proof}
Az előző állítás kiemelt azonosságát \emph{Bezout-azonosságnak}\index{Bezout-azonosság}
mondjuk.
}
\frame{
\begin{proposition}
	Legyenek a $p_1,\ldots,p_r\in\mathbb{F}\left[ t \right]$ tetszőleges $\mathbb{F}$ test feletti polinomok.
	Ezek pontosan akkor relatív prímek,
	ha léteznek $f_1,\ldots,f_r\in\mathbb{F}\left[ t \right]$ polinomok, hogy
	\[
		f_1\left( t \right)p_1\left( t \right)+f_2\left( t \right)p_2\left( t \right)+\ldots+f_r\left( t \right)p_r\left( t \right)=1\qedhere
	\]
\end{proposition}
}
\frame{
\begin{proposition}
	Legyenek a $p,q\in\mathbb{F}\left[ t \right]$ polinomok relatív prímek,
	és tegyük fel, hogy
	\begin{math}
		p|qr.
	\end{math}
	Ekkor $p|r$.
    \label{pr:rprim}
\end{proposition}
%\begin{proof}
%	Mivel a $p$ és a $q$ polinomok relatív prímek,
%	ezért a Bezout-azonosság\index{Bezout-azonosság} szerint van $f$ és $g$ polinom, amelyekre
%	\(
%	fpr+gqr=r.
%	\)
%	A feltétel szerint $qr$ a $p$ többszöröse,
%	így az iménti azonosság baloldala a $p$ többszöröse,
%	ami éppen azt jelenti, hogy $p|r$.
%\end{proof}
\begin{proposition}
	Tekintsük a $p,q\in\mathbb{F}\left[ t \right]$ normált polinomokat.
	Jelölje $d$ a legnagyobb közös osztót,
	és
	$m$ a legkisebb közös többszöröst.
	Ekkor
	\begin{displaymath}
		d\left( t \right)m\left( t \right)=p\left( t \right)q\left( t \right).\qedhere
	\end{displaymath}
\end{proposition}
%\begin{proof}
%	Legyen $p=dr_1$ és $q=dr_2$.
%	Először megmutatjuk, hogy $r_1$ és $r_2$ relatív prímek.
%	A Bezout-azonosság szerint valamely $f,g$ polinomra
%	\(
%	d=fdr_1+gdr_2.
%	\)
%	Kihasználva, hogy nullosztómentes gyűrűben nem zérus elemmel egyszerűsíteni lehet,
%	kapjuk az
%	\(
%	1=fr_1+gr_2
%	\)
%	azonosságot.
%	Ez azt jelenti, hogy $r_1$ és $r_2$ relatív prímek.
%	%    Ha $s$ polinom közös osztójuk, akkor $ds$ is közös osztója $p$-nek és $q$-nak,
%	%    amiből $ds|s$ következik.
%	%    A fokszámokat összehasonlítva ez csak akkor lehetséges, ha $s$ konstans polinom, 
%	%    azaz $s|1$ valóban fennáll.
%
%	Most megmutatjuk, hogy
%	\[
%		dr_1r_2\tag{\dag}
%	\]
%	a legkisebb közös többszörös.
%	Világos, hogy ez egy közös többszöröse a $p,q$ polinomoknak.
%	Tegyük fel, hogy $s$ egy másik közös többszörös, azaz
%	$s=ps_1$ és $s=qs_2$.
%	Ekkor
%	\begin{math}
%		dr_1s_1=ps_1=s=qs_2=dr_2s_2,
%	\end{math}
%	amiből újra a nullosztómentesség szerint
%	\[
%		r_1s_1=r_2s_2.
%	\]
%	Na most,
%	a fent kiemelt azonosság szerint szerint $r_1|r_2s_2$, ahol $r_1$ és $r_2$ relatív prímek.
%	Ebből azonnal kapjuk, hogy $r_1|s_2$.
%	No de, $s=qs_2=dr_2s_2$, amiből már látszik,
%	hogy $s$ egy többszöröse a $dr_1r_2$ polinomnak.
%	Az is világos, hogy (\dag) normált polinom, ez az $m$ legkisebb közös többszörös.
%	Innen
%	$d\cdot m=d(dr_1r_2)=(dr_1)(dr_2)=p\cdot q$ már nyilvánvaló.
%\end{proof}
}
\section{Az Euklideszi-algoritmus}
\frame{
Algoritmust keresünk polinomok legnagyobb közös osztójának és legkisebb közös többszörösének meghatározására.
%Ha egy pillanatra $\left( p_1,\ldots,p_n \right)$ jelöli az adott $p_1,\ldots,p_n$ polinomok legnagyobb közös osztóját,
%akkor nem nehéz meggondolni,
%hogy
%\[
%	\left( \left( p_1,\ldots,p_{n-1} \right),p_n \right)=\left( p_1,\ldots,p_{n}\right).
%\]
%Ezt $n=3,4,\ldots$ számokra alkalmazva azt kapjuk,
%hogy ha módszerünk van két polinom legnagyobb közös osztójának meghatározására,
%akkor evvel már akárhány -- persze véges sok -- polinom legnagyobb közös osztója is meghatározható.
%Analóg módon ugyanez igaz a legkisebb közös többszörösre is.
%Azt gondoltuk meg tehát, hogy ha meg tudnánk határozni két polinom legnagyobb közös többszörösét és legkisebb közös osztóját,
%akkor ugyanezt már meg tudnánk tenni több polinom esetén is.
%
%Az Euklideszi-algoritmus két polinom legnagyobb közös osztójának meghatározására szolgál.
%Az eddigi ismereteink szerint a $p,q$ legnagyobb közös osztója a $J\left( p,q \right)$ ideál legalacsonyabb
%fokú, normált tagja.
Az Euklideszi-algoritmus egy véges sok lépésben végrehajtható algoritmus,
egy a definíciónál sokkal hatékonyabb eljárás, két polinom generálta főideál generáló elemének,
azaz a legnagyobb közös osztónak a meghatározására.
}
\frame{

Az algoritmus megértése előtt emlékeznünk kell arra,
hogy ha $p,q\in\mathbb{F}\left[ t \right]$ valamely polinomok, akkor
\(
J\left( p,q \right)=\left\{ fp+gq:f,g\in\mathbb{F}\left[ t \right] \right\}
\)
jelöli a $p$ és a $q$ polinomokat tartalmazó legszűkebb ideált.

Az algoritmus lényege, hogy a $J\left( p,q \right)$ ideál két generátor elemét egyre alacsonyabb fokú
polinomokkal helyettesítjük, amíg megtaláljuk az ideál legalacsonyabb fokú polinomját.
\begin{lemma}
    Tegyük fel, hogy a $p,q,h,r\in\mathbb{F}[t]$ polinomokra
    \[
    p=hq+r.
    \]
    Ekkor a $p,q$ generálta $J\left( p,q \right)$ ideál egybeesik a $q,r$ generálta $J\left( q,r \right)$ ideállal, azaz 
    \[
        J\left( p,q \right)=J\left( q,r \right).
    \]
\end{lemma}
}
\frame{
\begin{proposition}[Euklidesz]\index{Euklideszi-algoritmus}
	Legyen $p,q\in\mathbb{F}\left[ t \right]$ nem zérus polinomok.
	Definiálja $p_{-1}=p,p_{0}=q$.
	Folytatva, ha valamely $i\geq 0$ számra $p_{i-1}$ és $p_i$ már definiált és $p_i\neq 0$,
	akkor a maradékos osztás szerint létezik egyetlen $h_{i+1},p_{i+1}\in\mathbb{F}\left[ t \right]$ polinom,
	amelyre
	\[
		p_{i-1}=h_{i+1}p_i+p_{i+1};
		\text{ ahol }
		\deg p_{i+1}<\deg p_i.\tag{\dag}
	\]
	Mivel minden egyes lépésben csökken a fokszám,
	ezért van olyan $s\geq 0$,
	hogy $p_s\neq 0$, de $p_{s+1}=0$.
	\\
	Erre a $p_s$ polinomra $J\left( p_s \right)=J\left( p,q \right)$, ezért
	$p_s$ normáltja a $p\text{ és a }q$ polinomok legnagyobb közös osztója.
\end{proposition}
%\begin{proof}
%	A fenti algoritmussal olyan
%	$p_{-1}, p_0,p_1,\ldots,p_s,p_{s+1}$ polinomokat kapunk,
%	amelyekre minden $i=0,\ldots,s$ mellett a (\dag) azonosság fennáll,
%	és $\deg p_{s+1}=-\infty$,
%	azaz $p_{s+1}=0$.
%
%	A (\dag) azonosság szerint,
%	minden $i=0,1,\ldots,s$ mellett
%	\begin{math}
%		p_{i-1}\in J\left( p_i,p_{i+1} \right)
%	\end{math},
%	amiből a
%	\begin{math}
%		J\left( p_{i-1},p_i \right)\subseteq J\left( p_i,p_{i+1} \right)
%	\end{math}
%	tartalmazás adódik.
%	Másrészt a (\dag) azonosságot értelmezhetjük úgy is, hogy
%	\(
%	p_{i+1}\in J\left( p_{i-1},p_i \right)
%	\),
%	amiből persze
%	\begin{math}
%		J\left( p_i,p_{i+1} \right)\subseteq J\left( p_{i-1},p_i \right)
%	\end{math}
%	következik.
%	Így minden $i=0,\ldots,s$-re végül is
%	\[
%		J\left( p_{i-1},p_i \right)
%		=
%		J\left( p_i,p_{i+1} \right).
%	\]
%	Alkalmazva ezt minden $i=0,\ldots,s$ mellett
%	\[
%		J\left( p,q \right)
%		=
%		J\left( p_{-1},p_0 \right)
%		=
%		J\left( p_0,p_1 \right)
%		=
%		J\left( p_1,p_2 \right)
%		=\ldots=
%		J\left( p_s,p_{s+1} \right)
%		=J\left( p_s \right).\qedhere
%	\]
%\end{proof}
}
\frame{
    Emlékezzünk arra, hogy két polinom szorzata a legangyobb közös osztó és a legkisebb közös többszörös szorzata:
    \[
        dm=pq.
    \]
Úgy interpretáljuk,
hogy ha a $pq$ szorzatot osztom maradékosan a legnagyobb közös osztóval, akkor a maradék mindig
zérus,
és a hányados éppen a legkisebb közös többszörös.

Azt gondoltuk meg tehát,
hogy az Euklideszi-algoritmus módszert ad két polinom legkisebb közös többszörösének algoritmikus meghatározására is.
}
\frame{
Ilyen módon véges sok lépésben végrehajtható algoritmust kapunk véges sok polinom legkisebb közös többszörösének meghatározására.

Például öt polinom legkisebb közös többszöröséhez,
egy Euklideszi-algoritmussal meghatározzuk az első kettő polinom legnagyobb közös osztóját,
majd egy újabb maradékos osztással az első kettő legkisebb közös többszörösét.
Ugyanezt teszem az így kapott és a harmadik polinommal,
az eredmény az első három polinom legkisebb közös többszöröse.
Az így kapott polinommal és a negyedik polinommal egy újabb Euklideszi-algoritmus és egy újabb maradékos osztás után kapjuk az első négy polinom legkisebb közös többszörösét,
majd ennek eredményével és az ötödik polinommal mint két polinomnak a legkisebb közös többszörösével kapjuk az eredeti öt polinom legkisebb közös többszörösét.
}

\end{document}
